% !TeX root = ../thesis.tex
% TH-CH03-WRITE-V1: stage 3 draft; not compiled or frozen.
\chapter{面向灵巧手操作的抓取候选生成与专家示教构建}
\label{chap:grasp_demonstrations}

\section{抓取任务定义与总体流程}
\label{sec:grasp_task_pipeline}

第二章建立了机器人执行与触觉观测的仿真基础，但具有可控模型并不意味着能够直接生成有效操作轨迹。面对不同形状和摆放状态的物体，系统既要寻找手部可形成接触的构型，也要使该构型满足机械臂运动约束，并在实际动力学推进后保持物体。因此，本章将专家示教构建组织为候选生成、机器人约束筛选、轨迹执行和质量验收四个层次，研究从物体几何输入到可保存控制序列的转换，而不将某一几何评分直接作为任务成功判据。

抓取输入包括物体标识、表面采样及法向、初始位姿与任务条件。候选以物体坐标系下的手部平移、旋转和六维驱动比例表示，可记为
\begin{equation}
 \mathcal G=(\boldsymbol t_{oh},\boldsymbol R_{oh},\boldsymbol f,\mathcal C,\boldsymbol m),
 \qquad \boldsymbol f\in[0,1]^6,
 \label{eq:ch03_candidate}
\end{equation}
其中，$\boldsymbol t_{oh}\in\mathbb R^3$和$\boldsymbol R_{oh}\in SO(3)$描述物体坐标系下的手部目标，$\mathcal C$存储接触点、法向及距离，$\boldsymbol m$为候选质量和来源指标。驱动比例需进一步转换为具体执行器的控制范围，不等于十二个手部关节的独立角度。候选合同检查维度、旋转合法性和比例范围，这些检查只能保证表达有效，不能证明实际接触已经成立。

\begin{figure}[htbp]
 \centering
 \fbox{\begin{minipage}{0.91\textwidth}\centering
 物体几何、初始位姿和任务条件\\[4pt]
 $\downarrow$\\[4pt]
 几何候选与物理反馈细化 $\rightarrow$ 机器人目标转换与可达筛选\\[4pt]
 $\downarrow$\\[4pt]
 局部腕姿态扩展与分段执行 $\rightarrow$ 轨迹记录\\[4pt]
 $\downarrow$\\[4pt]
 独立Strict质量检查 $\rightarrow$ 任务专属示教与结果记录
 \end{minipage}}
 \caption{抓取候选到专家示教的分级构建流程}
 \label{fig:ch03_overall_pipeline}
\end{figure}

图\ref{fig:ch03_overall_pipeline}中，机器人约束筛选将对象相对目标转换为世界坐标下的机械臂末端目标，轨迹执行再连接接近、闭合、保持和后续操作。只有经过执行，才能获得状态、接触和物体运动等反馈；Strict质量验收进一步检查尾段稳定性。Lift提供抓持和抬升的基本案例，搬移与堆叠可在抓持前缀上构造任务后缀，推动和装配则需要额外的接触或姿态控制。五类任务共享示教数据合同，但不共享完全相同的物理执行逻辑。本章重点讨论示教如何形成，记录关联、数据转换与全量质量审计留在第四章。

\section{多来源抓取候选生成}
\label{sec:grasp_candidate_generation}

欠驱动手的关节并非独立受控，直接优化全部关节角可能偏离实际可执行构型。本文以六个驱动比例为手部参数，利用仿真采样拟合的皮肤点轨迹代理连接驱动与接触几何。该代理用于可微评价，其范围受采样和模型版本限制，最终行为仍需物理反馈检验。

GraspQP通过可微抓取质量评价支持候选优化\cite{zurbrugg2025graspqp}。本项目调用其官方指标工厂，结合物体表面点、法向和手部代理构造接触评价。优化变量包括三维平移、六维旋转参数及六个驱动比例的logit，驱动比例通过sigmoid限制在有效区间。旋转参数在评价时恢复为旋转矩阵，从而避免直接优化九个任意矩阵元素。物体采样和手部皮肤点变换到同一坐标下后，程序按手指接触点组计算到物体的距离，并以软权重构造接触位置和法向。由此形成的接触集合用于官方质量指标计算，而不是把最近点距离本身当作力闭合证明。

单独优化抓取指标可能得到难以执行的构型，因此目标还加入接触距离、穿透、桌面距离、平移漂移和驱动耦合约束。接触距离项鼓励手部采样点接近物体；穿透和桌面项限制明显不合理的几何关系；漂移项抑制手部偏离初始种子过远；驱动耦合项约束部分手指比例的离散程度。各项具有不同归一化尺度与权重，其数值意义应结合实际实现解释。表\ref{tab:ch03_candidate_parameters}列出变量和约束的职责。本文在这一环节完成的是面向具体手部的指标适配与流程集成，不将GraspQP本身作为原创方法，也不声称几何目标已覆盖所有接触稳定性条件。

\begin{table}[htbp]
 \centering\small
 \caption{候选优化参数及约束的实际作用}
 \label{tab:ch03_candidate_parameters}
 \begin{tabular}{p{0.20\textwidth}p{0.25\textwidth}p{0.41\textwidth}}
 \hline
 项目 & 表达或输入 & 作用与边界\\
 \hline
 手部平移与旋转 & 三维平移、六维旋转参数 & 搜索对象相对手部位姿，输出合法旋转矩阵\\
 手部驱动 & 六维比例logit & 经sigmoid形成执行器比例，不独立控制全部关节\\
 接触几何 & 表面点、法向及距离 & 连接手部代理与物体几何评价\\
 官方抓取指标 & 接触集合与摩擦设置 & 提供几何抓取质量，不代替机器人执行\\
 附加几何约束 & 接触、穿透、桌面距离 & 排除部分明显不合理的候选\\
 漂移与驱动约束 & 种子平移及驱动比例 & 限制局部搜索偏移和驱动分散\\
 \hline
 \end{tabular}
\end{table}

几何候选进入DexEvolve-style细化阶段后，评价对象从接触构型扩展到执行行为。DexEvolve原研究采用进化优化获得多样且稳健的灵巧抓取\cite{zurbrugg2026dexevolve}，本项目保留候选群体、变异、交叉和档案管理的组织思想，并针对六驱动手及现有仿真执行接口作适配。平移、旋转和驱动比例分别施加扰动，驱动变化限制在有效范围内；交叉则组合不同候选的位姿与手部配置。候选在执行前还可进行接触修正和去穿透处理，使搜索不仅依赖最初的几何种子。

物理反馈通过实际执行候选获得接触保持、物体相对漂移和扰动响应等指标，必要时使用MJWarp后端评价寿命相关目标。几何质量较高的构型仍可能在接近时偏离物体，或在加载后因接触方向不合适而失稳，因而物理反馈不能由静态接触分数完全替代。进化档案利用候选距离、新颖性及局部密度等信息保留不同构型，避免全部预算集中到近乎重复的局部解。但档案多样性是搜索机制，不直接等于独立初态或不同成功模式已经得到充分验证。本项目使用的是适配实现，并未据此声称完整复现第三方Isaac系统。

\begin{figure}[htbp]
 \centering
 \fbox{\begin{minipage}{0.91\textwidth}\centering
 物体表面与六驱动手部代理\\[4pt]
 $\downarrow$\\[4pt]
 GraspQP官方指标 $+$ 本地几何约束 $\rightarrow$ 初始优化候选\\[4pt]
 $\downarrow$\\[4pt]
 DexEvolve-style变异、交叉、接触修正与物理反馈\\[4pt]
 $\downarrow$\\[4pt]
 候选档案 $\rightarrow$ Wrist Lattice局部扩展\\[4pt]
 $\downarrow$ 未解决且满足后续条件\\[4pt]
 MJWarp PPO模板选择与驱动编辑
 \end{minipage}}
 \caption{已有候选模块在分级生成流程中的职责}
 \label{fig:ch03_candidate_sources}
\end{figure}

Wrist Lattice针对候选与具体机械臂的可达关系开展局部扩展。它在源候选附近改变腕部平移和朝向，再以机器人逆运动学残差排序，保留可进一步执行的模板。当前默认集合包含39个格点，平移基本尺度为0.01~m，旋转基本尺度为$15^\circ$，并包含若干倍尺度的单轴旋转以及平移、旋转耦合编辑。局部改变不会把手部根节点整体绕物体远距离扫动，适合在已有接触构型附近寻找机械臂更易跟踪的目标。其局限是搜索依赖源候选和固定邻域，不能保证发现相距较远的可达抓取模式；39个格点也不是39个独立实验样本。

统一生成预算配置当前采用64个种子、150步GraspQP优化并执行12个候选；进化阶段使用24个初始个体、每代12个子代、16代，名义候选评价次数为216，随后向腕部扩展提供最多6个档案候选。这些是当前入口默认预算；历史批次及缓存复用的实际代价需另行区分。

图\ref{fig:ch03_candidate_sources}中的MJWarp PPO仅作为条件性补充，在格点阶段未解决且仍有可用模板时进行模板选择与手部驱动编辑。历史Top100路由中存在2条RL\_SUCCESS记录，但这并不构成强化学习优于其他模块的消融证据。整体而言，可微指标提供几何方向，进化反馈检验物理行为，格点扩展处理机器人可达性，后续编辑尝试解决局部未完成案例。各模块依赖不同输入和预算，应分别说明作用，不能混合命名为一个新的原创抓取优化算法。

\section{基于可达性约束的抓取筛选}
\label{sec:grasp_reachability}

对象相对抓取只有在转换到机器人末端目标后，才能评价机械臂是否能够执行。设物体世界位姿为$(\boldsymbol R_{wo},\boldsymbol p_{wo})$，候选定义的手部目标为$(\boldsymbol R_{oh},\boldsymbol t_{oh})$，则目标手部姿态与末端姿态满足
\begin{equation}
 \begin{aligned}
 \boldsymbol p_{wh}&=\boldsymbol p_{wo}+\boldsymbol R_{wo}\boldsymbol t_{oh},\\
 \boldsymbol R_{wh}&=\boldsymbol R_{wo}\boldsymbol R_{oh},\qquad
 \boldsymbol R_{we}=\boldsymbol R_{wh}\boldsymbol R_{eh}^{\mathsf T}.
 \end{aligned}
 \label{eq:ch03_target_transform}
\end{equation}
其中，$\boldsymbol R_{eh}$为手部安装相对旋转，$w,o,h,e$分别表示世界、物体、候选手部参考与末端坐标。位置量单位为m，旋转无量纲。这里采用执行器接口中实际使用的安装参考，不能把STL局部原点任意替换为另一个掌心点后仍沿用同一公式。若缺少安装补偿，即使目标手部方向在几何上正确，机械臂实际跟踪方向也可能与之不同。

预抓取目标沿候选手部局部方向向外退让，当前默认距离为0.09~m；抓取目标则对应候选位置本身。筛选器按预抓取、抓取的顺序调用实际机械臂IK，在前一个求解结果上继续处理后一个目标。它同时记录位置残差、旋转残差与末端行程，用于区别几何合理但机器人难以到达的候选。检查过程中会保存并恢复环境状态、执行器控制量以及控制器历史目标，避免静态候选排序直接改变后续真正执行的起点。

当前排序器采用如下评分：
\begin{equation}
 S_{\mathrm{reach}}=20e_p^{\max}+0.6e_R^{\max}+0.08\ell,
 \label{eq:ch03_reachability_score}
\end{equation}
其中，$e_p^{\max}$和$e_R^{\max}$分别为两个目标上的最大位置与旋转残差，$\ell$为相应末端行程。三个系数是代码中的数值权重，用于产生排序，而非经过物理标定的统一距离单位。评分较小表示在该初值和求解配置下更易到达，不表示对应抓持必定稳定。静态排序还临时放宽目标速度约束，使残差更接近求解可达性的检查；真正执行时仍使用相应速度限制和平滑，不能将两条路径的跟踪结果视为相同。

模板构建依据执行配置保留位置误差不超过0.025~m、旋转误差不超过0.22~rad的候选。这两个容差属于当前生成执行层的筛选设置，不是第二章控制器内部更严格的迭代终止容差。满足筛选容差后，还需要根据腕部编辑幅度、源轨迹动态跟踪误差和候选评分确定执行顺序。这样既考虑静态残差，也保留对原有接触构型的偏好，避免搜索一开始就使用偏移很大的变体。

\begin{figure}[htbp]
 \centering
 \fbox{\begin{minipage}{0.91\textwidth}\centering
 源候选的局部腕姿态编辑：平移、旋转及耦合格点\\[4pt]
 $\downarrow$\\[4pt]
 对象相对目标 $\rightarrow$ 世界手部目标 $\rightarrow$ 安装旋转补偿\\[4pt]
 $\downarrow$\\[4pt]
 预抓取目标IK $\rightarrow$ 抓取目标IK\\[4pt]
 $\downarrow$\\[4pt]
 残差阈值筛选 $\rightarrow$ 评分、编辑幅度与源跟踪误差排序\\[4pt]
 $\downarrow$\\[4pt]
 有限预算下的动态执行；不是全轨迹无碰撞认证
 \end{minipage}}
 \caption{局部腕姿态扩展与两个目标的可达性筛选}
 \label{fig:ch03_reachability}
\end{figure}

图\ref{fig:ch03_reachability}所示预检的作用范围需要明确：已核验实现主要覆盖预抓取与抓取两个目标，并没有据此证明完整Lift轨迹的每个路点都通过连续碰撞认证。两个端点均可达时，中间路径仍可能靠近桌面或改变IK分支；静态残差较小时，动态跟踪也可能因速度限制而产生更大误差。因此，筛选结果只决定候选是否值得进入执行，后续阶段仍需读取真实状态和接触结果，不能把它称为通用连续无碰撞规划器。

不同约束在流程中承担不同职责。候选优化中的桌面距离惩罚限制明显不合理的初始几何，任务迁移中的保守下降高度控制局部接近，执行过程监测接触和目标误差，Strict则对尾部桌面接触及稳定性作最终判断。这些机制相互补充，但不应合并为一个未经实现的全局碰撞约束。对失败也应保留层次：无有效候选、静态IK残差过大、动态接近偏差、闭合后缺少对向接触和抬升失稳对应不同问题。只有明确失败发生在哪一层，才能避免把所有失败归因于抓取几何质量。

局部搜索还受预算限制。模板构建对可执行候选排序后，只运行有限数量的完整尝试，并可优先复用已有参考。这有助于控制批量生成代价，但意味着最终保留下来的模板不等于全部可达候选，也不构成对整个抓取空间的穷举。本文将这种筛选视为面向具体机器人和源候选的约束连接环节，其有效范围由后续实际执行与验收结果共同确定。

\section{连续逆运动学与轨迹验证}
\label{sec:grasp_trajectory_verification}

通过静态筛选后，候选还需要变成从机器人当前状态出发的连续命令。执行器将目标分解为稳定初态、过渡、预抓取、接近、闭合、保持、抬升和验证等阶段，每个阶段具有目标末端位姿、手部命令和控制步数。分段组织的目的不是人为把操作拆成互不相关的动作，而是在接触建立前后采用不同目标，并保留阶段边界上的误差及接触信息。前一阶段输出的实际状态构成后一阶段的起点，因而不能只保存各阶段名义目标而忽略执行结果。

对单段执行，位置和手部控制量在起点与终点之间插值，姿态采用四元数SLERP。当启用平滑选项时，归一化进度经过三次映射后再用于插值：
\begin{equation}
 \begin{aligned}
 \bar u&=3u^2-2u^3,\qquad u\in[0,1],\\
 \boldsymbol p(u)&=(1-\bar u)\boldsymbol p_0+\bar u\boldsymbol p_1,\\
 \boldsymbol a_h(u)&=(1-\bar u)\boldsymbol a_{h,0}+\bar u\boldsymbol a_{h,1}.
 \end{aligned}
 \label{eq:ch03_segment_interpolation}
\end{equation}
其中，$\boldsymbol a_h$为六个手部执行器的位置命令；关闭平滑时直接使用$u$。该映射使参数曲线端点斜率为零，能够缓和相邻阶段的名义速度突变，但并不保证实际物体速度也为零。接触力、手部加载和关节伺服仍影响动力学结果，不能把插值平滑等同于抓持稳定。

末端目标按控制周期输入第二章所述IK控制器，控制器以已有目标和状态作为后续求解的参考，并限制关节更新与目标速度。这里不重复DLS推导，而强调其在连续执行中的角色：目标位姿是上层轨迹要求，实际关节状态由受限求解、执行器伺服和物理推进共同产生。执行过程保存接近位置和旋转残差，若超出执行配置允许范围，则停止将该候选视为有效接近。对接触建立后的阶段，还需要检查拇指对向皮肤接触、物体抬升和保持状态，防止仅凭完成命令序列就认定任务成功。

\begin{figure}[htbp]
 \centering
 \fbox{\begin{minipage}{0.91\textwidth}\centering
 当前真实状态 $\rightarrow$ 分段插值目标 $\rightarrow$ IK与位置命令\\[4pt]
 $\downarrow$ 物理推进与阶段记录\\[4pt]
 完整轨迹 $\rightarrow$ 按低层ctrl独立回放\\[4pt]
 $\downarrow$\\[4pt]
 尾段抬升、速度、角速度、皮肤接触、桌面接触及穿透\\[4pt]
 $\downarrow$ 联合判定\\[4pt]
 FINAL\_VERIFIED / FINAL\_REJECTED
 \end{minipage}}
 \caption{分段轨迹执行与Strict Replay质量验收}
 \label{fig:ch03_strict_replay}
\end{figure}

已有普通Replay接口主要检查控制序列能否回放，以及尾段任务成功比例。其成功条件要求最后一个尾段标记为成功，且尾段成功比例至少为0.8。该接口并未完整表达严格示教所需的稳定性条件，因此另外采用图\ref{fig:ch03_strict_replay}所示的Strict Replay。在对应仿真环境中重放实际低层控制序列后，Strict同时检查抬升、速度、角速度、皮肤接触和环境碰撞等指标，当前final配置见表\ref{tab:ch03_strict_thresholds}。

\begin{table}[htbp]
 \centering\small
 \caption{当前final配置的Strict Replay验收条件}
 \label{tab:ch03_strict_thresholds}
 \begin{tabular}{p{0.39\textwidth}p{0.17\textwidth}p{0.30\textwidth}}
 \hline
 指标 & 阈值 & 判定范围\\
 \hline
 最终及尾段最小抬升 & $\geq0.055$~m & 二者均需满足\\
 物体线速度 & $\leq0.10$~m/s & 尾段最大值\\
 物体角速度 & $\leq0.10$~rad/s & 尾段最大值\\
 皮肤接触、抓持比例 & 均$\geq0.80$ & 尾段统计\\
 机器人与桌面接触比例 & $\leq0$ & 尾段统计\\
 最大穿透 & $\leq0.003$~m & 尾段最大值\\
 验证尾段、额外保持 & 20步、12步 & 当前配置参数\\
 \hline
 \end{tabular}
\end{table}

Strict中的皮肤接触与抓持具有具体语义，不将任意机器人连杆碰撞当作有效指面抓持。全程最大穿透和尾段最大穿透虽然都被记录，实际最终布尔条件使用后者；若将全程记录量误写为硬门槛，就会改变原验证协议。最终及尾段最小抬升同时受限，也意味着物体曾短暂达到目标高度并不足够，尾段仍需满足要求。速度和角速度限制进一步区分持续承载与明显滑移、旋转等不稳定状态，桌面接触条件则避免利用环境支撑冒充稳定悬持。

验收状态由实际profile及布尔结果共同确定，final配置通过时记为FINAL\_VERIFIED，否则为FINAL\_REJECTED。生成过程中的某个变量即使名为strict\_execution，也不能替代该最终状态。类似地，物理执行返回success与严格回放通过是不同证据，候选档案、模板清单和示教导出需要保留验证层级。导出入口默认先进行final检查，诊断性允许未验证轨迹的开关不能作为正式质量保证。

轨迹重放还必须使用正确的动作语义。旧生成记录中的action可以是末端IK指令，ctrl则保存实际施加的机械臂七通道与手部六通道位置命令。二者均可能为十三维，但直接用前者替代后者会改变回放问题。记录器保存的关节状态、对象姿态和控制量有助于检查这种对应，完整时间关联则在第四章讨论。Strict通过只说明给定模型、初态、控制序列及阈值下的结果，不认证任意后续模型版本均能复现，也不把同一轨迹的后缀当成新的独立专家示教。

\section{多任务专家示教生成}
\label{sec:multitask_demonstration}

多任务示教共享候选来源、控制序列、对象状态和阶段信息等记录合同，但采用不同控制逻辑。Lift提供基本抓持案例，其他任务依据操作目标构造后续过程，并以实际物体状态而非执行器目标判断结果。主要阶段见表\ref{tab:ch03_task_stages}。

Lift执行从稳定初态开始，经过过渡和预抓取接近，在close阶段逐步建立手部闭合命令，在hold阶段检查和维持接触，然后抬升并进入verify。接近阶段需要机械臂跟踪手部目标，接触阶段则要求驱动命令与物体几何形成有效加载；两者失败不能用同一几何指标解释。完整阶段记录让后续能够判断失败发生在物体接触之前还是加载之后，也为其他任务选择可用抓持前缀提供依据。

PickPlace以源抓持为基础，增加抬高避让、横向运输、下降、释放、撤离及任务验证。转移时既要保留物体初态与抓持目标的关系，也要考虑手部和桌面之间的几何余量。当前实现存在恢复已验证Lift终态并复用原前缀的路径，也存在从reset重新执行的选择。前者可以避免重复执行中的数值漂移，却不能被描述成整个拼接序列在该次运行中全部重新产生。对示教而言，前缀来源及衔接方式应明确保留，不能只凭拼接后数组长度判断物理连续性。

后续在线记录工作对原有前缀复用分支作了区分，重新执行时保留实际观测和控制，不将历史Lift触觉直接复制成新采集结果。这项修复说明示教生成与数据保存之间存在接口约束：任务控制完成和记录正确并不是同一个条件。本节只交代两种生成路径及其限制，完整模态关联与修复验收由第四章展开，不把修复后的合同回填为全部历史文件都已具备的属性。

\begin{table}[htbp]
 \centering\small
 \caption{五类任务的示教阶段与主要执行约束}
 \label{tab:ch03_task_stages}
 \begin{tabular}{p{0.15\textwidth}p{0.42\textwidth}p{0.29\textwidth}}
 \hline
 任务 & 主要阶段 & 执行重点\\
 \hline
 Lift & settle、transit、pregrasp、approach、close、hold、lift、verify & 接近跟踪、对向接触与抬升保持\\
 PickPlace & 抓持前缀、clear、transport、descend、release、retreat、task\_verify & 来源保持、安全下降及释放后目标状态\\
 Stack & 抓取、目标定位、运输下降、释放、稳定性评价 & 被操作物与支撑物关系\\
 Push & 接触接近、接触检测、反馈推动、目标评价 & 持续推动与物体实际位移\\
 NutAssembly & Bolt抓持、姿态伺服、平移、下降、装配评价 & 载荷保持、朝向与装配几何\\
 \hline
 \end{tabular}
\end{table}

Stack除被操作物之外还需要支撑物的状态。源抓持能够提供手与物体的相对关系，但堆叠目标必须根据支撑物位置、物体尺寸和下降范围进一步确定。任务后缀通过运输、下降和释放形成支撑关系，再检查物体是否稳定放置。并非每个物体都适合当前堆叠定义，因此已有批次先选取task-valid对象域。对这一筛选应保留来源，不能将限定物体域上的成功比例写成对任意物体的覆盖能力。

\begin{figure}[htbp]
 \centering
 \fbox{\begin{minipage}{0.91\textwidth}\centering
 Lift抓持与承载记录 $\rightarrow$ PickPlace搬移释放后缀\\[4pt]
 \phantom{Lift抓持与承载记录} $\rightarrow$ Stack支撑与稳定放置后缀\\[5pt]
 Push：几何接近 $\rightarrow$ 有效接触 $\rightarrow$ 反馈推动\\[5pt]
 NutAssembly：Bolt抓持 $\rightarrow$ 姿态伺服 $\rightarrow$ 对准下降\\[5pt]
 共享DemonstrationEpisode合同，不共享完全相同的控制逻辑
 \end{minipage}}
 \caption{多任务示教的共用来源与任务专属控制关系}
 \label{fig:ch03_task_transfer}
\end{figure}

Push不以先形成稳定悬持为必要前提，而是依据手掌或手指几何建立推动接触，检测物体接触并按目标状态反馈推进。手部到达目标附近并不说明物体已经产生期望位移，因此需要单独检查实际运动与接触条件。该任务不能直接继承Lift的成功定义，触觉激活也只能作为接触相关信息，不能替代目标位移评价。图\ref{fig:ch03_task_transfer}因此将推动作为专属控制链，而非强行接在抬升轨迹之后。

NutAssembly包含抓持Bolt、姿态伺服、平移、下降与装配结果评价。抓持阶段的细小偏差可能使物体倾斜，在后续下降前就丢失可用接触，因此不能把所有装配失败都归因于对孔精度。现有实现保留手部控制策略对承载过程的影响，也保留旧variant与后续几何版本不兼容的记录。局部成功可以用于分析某个控制条件是否值得保留，却不足以说明不同初态都能稳定装配，更不能把旧版本成功清单当成当前几何的最终验收。

上述专家示教主要由候选、规划和任务反馈控制形成，与人工遥操作获得示教的方式不同\cite{mandlekar2022demonstrations,mandlekar2023mimicgen}。其价值在于能明确生成条件和执行判据，但仍需防止把名义目标当作已实现状态、把保存接口当作质量认证。示教轨迹进入后续学习或导出之前，应携带任务与验证信息；LeRobot转换、执行前后时序、lineage及去重等问题在第四章集中讨论。本章不以数据文件数量代替任务验证，也不将自动生成的专家轨迹等同于已经训练出可靠的闭环策略。

\section{抓取及示教质量实验}
\label{sec:grasp_quality_experiments}

本节区分历史批量路由状态与具有源manifest、数组及Strict结果的具体轨迹。文件核对不等于重新运行验证；候选成功、任务成功和FINAL通过具有不同分母，以下不将其合并，也不对不同条件的批次进行算法排名。

历史批量记录见表\ref{tab:ch03_historical_batches}。Lift候选路由涉及100个对象标识，其中98条登记为LATTICE\_SUCCESS，2条为RL\_SUCCESS，说明该批次主要通过格点模板阶段获得可用结果，不能据此推断进化优化或强化学习的独立贡献。PickPlace与Push汇总虽然均登记100条SUCCESS，但前者全部为cached，后者有98条cached，因此不能解释为本轮重新完成的两百次物理执行。Stack在task-valid筛选域中得到43条成功与5条失败，而NutAssembly的100条记录以variant为单位，不等于100种独立操作物体。

\begin{table}[htbp]
 \centering\small
 \caption{历史批量路由与任务记录及其统计限制}
 \label{tab:ch03_historical_batches}
 \begin{tabular}{p{0.16\textwidth}p{0.12\textwidth}p{0.30\textwidth}p{0.28\textwidth}}
 \hline
 任务或阶段 & 记录数 & 历史状态 & 解释限制\\
 \hline
 Lift候选路由 & 100 & 98 LATTICE\_SUCCESS，2 RL\_SUCCESS & 不是100条FINAL\_VERIFIED\\
 PickPlace & 100 & 100 SUCCESS & 全部cached\\
 Stack & 48 & 43 SUCCESS，5 FAILED & 仅task-valid对象域\\
 Push & 100 & 100 SUCCESS & 98条cached\\
 NutAssembly & 100 & 100 SUCCESS & variant记录，非100个独立物体\\
 \hline
 \end{tabular}
\end{table}

批次版本也影响这些记录的适用范围。Top100汇总保存18项源码Hash，本次材料核对发现12项与当前相同、5项不同、1项对应文件缺失。虽然其中手部XML及部分候选模块一致，仍不能把当前全部执行逻辑回填成历史运行配置。缓存行中的零耗时也不代表候选优化不需要计算。该组结果可说明分级流程已产生多任务历史产物，却不能作为统一模型、统一初态与统一Strict合同下的综合通过率。

局部扰动示教采用Sugar Box的10条既有记录。每份均有对应manifest，包含不同seed及不同NPZ Hash，源文件Hash与原审计登记一致，10份历史Strict结果均为FINAL\_VERIFIED。表\ref{tab:ch03_sugar_strict}列出其最终抬升和尾段角速度，展示严格通过并不只依赖数组中的成功标签。各记录使用同一物体和相近源抓取，历史审计记载不同初始位置的两两距离为0.707--2.828~mm，属于局部条件变化，不是大范围工作空间或跨物体泛化测试。

\begin{table}[htbp]
 \centering\small
 \caption{Sugar Box局部扰动示教的历史Strict结果}
 \label{tab:ch03_sugar_strict}
 \begin{tabular}{rrrrl}
 \hline
 序号 & seed & 最终抬升（m） & 尾段最大角速度（rad/s） & 验收状态\\
 \hline
 1 & 18011 & 0.06431 & 0.04049 & FINAL\_VERIFIED\\
 2 & 18012 & 0.06304 & 0.00324 & FINAL\_VERIFIED\\
 3 & 18013 & 0.06304 & 0.00303 & FINAL\_VERIFIED\\
 4 & 18014 & 0.06299 & 0.00311 & FINAL\_VERIFIED\\
 5 & 18015 & 0.06115 & 0.05234 & FINAL\_VERIFIED\\
 6 & 18016 & 0.06387 & 0.04748 & FINAL\_VERIFIED\\
 7 & 18017 & 0.06265 & 0.00602 & FINAL\_VERIFIED\\
 8 & 18018 & 0.06420 & 0.00428 & FINAL\_VERIFIED\\
 9 & 18019 & 0.06400 & 0.04943 & FINAL\_VERIFIED\\
 10 & 18020 & 0.06304 & 0.00346 & FINAL\_VERIFIED\\
 \hline
 \end{tabular}
\end{table}

这些记录支持源模板附近存在多个历史Strict通过序列。seed和Hash提供可追溯性，却不能消除同物体、同模板关联；10/10不是对未见场景的通过率估计，也不提供候选算法的公平对照。

\begin{table}[htbp]
 \centering\small
 \caption{历史90次扰动尝试的汇总分类}
 \label{tab:ch03_generation_failures}
 \begin{tabular}{lrp{0.51\textwidth}}
 \hline
 汇总类别 & 次数 & 统计解释\\
 \hline
 接受 & 46 & 历史生成汇总中的接受结果\\
 Strict不稳定 & 18 & 进入严格检查后未满足稳定条件\\
 Lift失败 & 26 & 未形成相应成功抬升轨迹\\
 \hline
 \end{tabular}
\end{table}

更大一组历史扰动生成报告登记90次尝试，其中46条接受、18条Strict不稳定、26条Lift失败，见表\ref{tab:ch03_generation_failures}。这组数据在本章保留为L1历史汇总，没有在本轮逐一重新认证全部90条原始轨迹。报告中的低抬升和角速度超限诊断还可能重叠，不应把所有诊断项当作互斥失败类别相加。该批次与上述Sugar示教存在来源重用，二者也不能累计成一百条新的独立验证。这些结果表明，能够生成几何或可达候选之后，承载与尾段稳定性仍需分别处理。

典型Lift案例的历史Strict结果中，最终抬升为0.0608926~m，尾段接触和抓持比例均为1，支持在该条件下持续承载的解释。NutAssembly则保留原始执行失败、降低抓点后仍失败、在降低抓点基础上保留hold抓持命令后单例成功的对照。该结果提示问题可能在装配之前的抓持保持阶段已经形成，而非只与末端对准有关；但它只是局部控制条件的对照，不能解释为跨seed稳定提升。图\ref{fig:ch03_examples}保留真实案例图位，后续只能从对应视频和manifest取得画面，不使用合成图片代替物理执行结果。

\begin{figure}[htbp]
 \centering
 \fbox{\parbox[c][2.0cm][c]{0.89\textwidth}{\centering
 真实案例图位：Lift历史Strict通过轨迹与Nut抓持保持对照。\\
 待从已有视频中提取并绑定对应manifest和阶段；\\
 当前不以示意性画面代替实验截图。}}
 \caption{典型成功与失败案例的待补原始画面}
 \label{fig:ch03_examples}
\end{figure}

现有结果支持分层评价与任务专属控制的必要性，尚不支持模块间的公平优越性结论。缓存、variant和局部成功不合并为统一成功率；未定位原件的14物体与127物体批次不列入验证结果，也不用其他数量替代。

\section{本章小结}
\label{sec:grasp_demonstration_summary}

本章围绕几何抓取候选如何转化为可执行专家示教，构建并说明了候选表达、几何与物理反馈细化、机器人约束筛选、连续轨迹执行及严格质量验收的工作链。GraspQP官方指标被适配到六驱动手部代理，DexEvolve-style搜索与局部腕姿态扩展分别补充执行反馈和机器人可达性信息，条件性后续编辑用于处理部分尚未解决的案例。这些环节具有不同作用，不以统一算法名称掩盖已有方法来源。

轨迹生成通过阶段目标、插值和控制器连续状态连接操作过程，Strict进一步区分任务成功与可用示教质量。已有批量记录、Sugar Box局部扰动Strict结果及典型任务对照提供了有限范围内的验证依据，但缓存记录、物体筛选域、历史版本及初态覆盖仍限制结论外推。五种任务共享记录合同，同时保留各自的接触、承载和目标评价逻辑。

专家示教生成只解决了控制序列从何而来的问题。轨迹进入后续学习和数据生产前，还需要检查多模态观测与动作的时间关系、执行来源和质量信息能否可靠关联。下一章据此讨论数据组织与质量审计，而不把示教文件的存在直接等同于训练数据已经完成验收。
